\documentclass[10pt]{extarticle}
\usepackage[a4paper,margin=1.5cm]{geometry}
\usepackage[T1]{fontenc}
\usepackage{amsmath,amssymb}
\usepackage{graphicx}
\usepackage[export]{adjustbox}
\usepackage{multicol}
\usepackage{enumitem}
\usepackage{booktabs,array,tabularx}
\setlength{\columnsep}{0.7cm}
\setlength{\parindent}{0pt}
\setlength{\parskip}{3pt}
\pagestyle{empty}
\makeatletter
\newcommand{\qfigdef}[4]{\@namedef{qf@file@#1}{#2}\@namedef{qf@w@#1}{#3}\@namedef{qf@wide@#1}{#4}}
\newcommand{\qf@check}[1]{\@ifundefined{qf@file@#1}{\PackageError{score2}{Unknown figure id #1}{}}{}}
\newcommand{\qf@img}[4][]{\includegraphics[width=#2\dimexpr\@nameuse{qf@w@#4}pt\relax,max width=#3,max height=0.45\textheight,#1]{\@nameuse{qf@file@#4}}}
\newif\ifqf@full
\newif\ifqf@boxed
\newsavebox{\qf@box}
\newcommand{\qf@open}{\global\qf@fulltrue\par\global\qf@boxedtrue
  \setbox\qf@box\vbox\bgroup\hsize\textwidth\linewidth\textwidth}
\newcommand{\qf@emit}[1]{\par\egroup\global\qf@boxedfalse
  \ifdim\dimexpr\ht\qf@box+\dp\qf@box\relax<#1\textheight\noindent\usebox{\qf@box}\par\else\unvbox\qf@box\fi}
\newcommand{\qfullsplit}{\ifqf@boxed\qf@emit{0.85}\fi}
\newcommand{\qf@close}{\ifqf@boxed\qf@emit{0.45}\fi\global\qf@fullfalse}
\newcommand{\qfullbegin}{\end{multicols}\qf@open}
\newcommand{\qfullend}{\qf@close\begin{multicols}{2}}
\newcommand{\qfullstart}{\qf@open}
\newcommand{\qfulldone}{\qf@close}
\newcommand{\qfig}[2][1]{\qf@check{#2}%
  \ifqf@full\qf@blk{#1}{#2}\else
  \ifnum\@nameuse{qf@wide@#2}=1 \expandafter\@firstoftwo\else\expandafter\@secondoftwo\fi
  {\end{multicols}\begin{center}\qf@img{#1}{\textwidth}{#2}\end{center}\begin{multicols}{2}}%
  {\qf@blk{#1}{#2}}\fi}
\newcommand{\qf@blk}[2]{\par\smallskip{\centering\qf@img{#1}{\linewidth}{#2}\par}\smallskip}
\newcommand{\qfigi}[2][1]{\qf@check{#2}\qf@img[valign=c]{#1}{\linewidth}{#2}}
\makeatother
\newcommand{\papersec}[1]{\par\bigskip{\centering\Large\bfseries PAPER-#1\par}\medskip}
\newcommand{\partsec}[2]{\par\medskip\noindent{\large\bfseries PART-#1: #2}\par\smallskip}
\newcommand{\sect}[1]{\par\smallskip\noindent\textbf{#1}\par}
\newcommand{\qn}[1]{\par\smallskip\noindent\textbf{#1.}~}
\newcommand{\paraq}[2]{\par\smallskip\noindent\textbf{Paragraph for Questions #1 and #2}\par}
\newenvironment{qoptions}{\begin{itemize}[leftmargin=1.8em,itemsep=1pt,topsep=2pt,label={}]}{\end{itemize}}
\newcommand{\opt}[2]{\item[(#1)] #2}
% figure registry (generated by score2_tools.py finalize)
\qfigdef{p1_phy_s1_q01_A}{SRG (Enthusiast) - Test 27 (13 Aug 2023)-img/p1_phy_s1_q01_A.png}{160.7}{0}
\qfigdef{p1_phy_s1_q04_1}{SRG (Enthusiast) - Test 27 (13 Aug 2023)-img/p1_phy_s1_q04_1.png}{95.6}{0}
\qfigdef{p1_phy_s1_q04_A}{SRG (Enthusiast) - Test 27 (13 Aug 2023)-img/p1_phy_s1_q04_A.png}{65.9}{0}
\qfigdef{p1_phy_s1_q04_B}{SRG (Enthusiast) - Test 27 (13 Aug 2023)-img/p1_phy_s1_q04_B.png}{62.4}{0}
\qfigdef{p1_phy_s1_q04_C}{SRG (Enthusiast) - Test 27 (13 Aug 2023)-img/p1_phy_s1_q04_C.png}{61.3}{0}
\qfigdef{p1_phy_s1_q04_D}{SRG (Enthusiast) - Test 27 (13 Aug 2023)-img/p1_phy_s1_q04_D.png}{61.7}{0}
\qfigdef{p1_phy_s1_q06_1}{SRG (Enthusiast) - Test 27 (13 Aug 2023)-img/p1_phy_s1_q06_1.png}{81.8}{0}
\qfigdef{p1_phy_s1_q07_1}{SRG (Enthusiast) - Test 27 (13 Aug 2023)-img/p1_phy_s1_q07_1.png}{98.5}{0}
\qfigdef{p1_phy_s1_q09_1}{SRG (Enthusiast) - Test 27 (13 Aug 2023)-img/p1_phy_s1_q09_1.png}{59.6}{0}
\qfigdef{p1_phy_s1_q11_1}{SRG (Enthusiast) - Test 27 (13 Aug 2023)-img/p1_phy_s1_q11_1.png}{98.4}{0}
\qfigdef{p1_phy_s1_q14_1}{SRG (Enthusiast) - Test 27 (13 Aug 2023)-img/p1_phy_s1_q14_1.png}{48.0}{0}
\qfigdef{p1_phy_s1_q15_1}{SRG (Enthusiast) - Test 27 (13 Aug 2023)-img/p1_phy_s1_q15_1.png}{69.4}{0}
\qfigdef{p1_phy_s1_q16_1}{SRG (Enthusiast) - Test 27 (13 Aug 2023)-img/p1_phy_s1_q16_1.png}{83.2}{0}
\qfigdef{p1_phy_s2_q01_1}{SRG (Enthusiast) - Test 27 (13 Aug 2023)-img/p1_phy_s2_q01_1.png}{136.4}{0}
\qfigdef{p1_phy_s2_q02_1}{SRG (Enthusiast) - Test 27 (13 Aug 2023)-img/p1_phy_s2_q02_1.png}{55.4}{0}
\qfigdef{p1_phy_s2_q03_1}{SRG (Enthusiast) - Test 27 (13 Aug 2023)-img/p1_phy_s2_q03_1.png}{104.4}{0}
\qfigdef{p1_phy_s2_q04_1}{SRG (Enthusiast) - Test 27 (13 Aug 2023)-img/p1_phy_s2_q04_1.png}{130.6}{0}
\qfigdef{p1_phy_s2_q05_1}{SRG (Enthusiast) - Test 27 (13 Aug 2023)-img/p1_phy_s2_q05_1.png}{79.4}{0}
\qfigdef{p1_phy_s2_q06_1}{SRG (Enthusiast) - Test 27 (13 Aug 2023)-img/p1_phy_s2_q06_1.png}{79.9}{0}
\qfigdef{p1_phy_s2_q07_1}{SRG (Enthusiast) - Test 27 (13 Aug 2023)-img/p1_phy_s2_q07_1.png}{69.3}{0}
\qfigdef{p1_phy_s2_q08_1}{SRG (Enthusiast) - Test 27 (13 Aug 2023)-img/p1_phy_s2_q08_1.png}{132.0}{0}
\qfigdef{p1_phy_s2_q09_1}{SRG (Enthusiast) - Test 27 (13 Aug 2023)-img/p1_phy_s2_q09_1.png}{62.9}{0}
\qfigdef{p1_phy_s2_q10_1}{SRG (Enthusiast) - Test 27 (13 Aug 2023)-img/p1_phy_s2_q10_1.png}{127.3}{0}
\qfigdef{p1_chem_s1_q02_A}{SRG (Enthusiast) - Test 27 (13 Aug 2023)-img/p1_chem_s1_q02_A.png}{65.4}{0}
\qfigdef{p1_chem_s1_q02_B}{SRG (Enthusiast) - Test 27 (13 Aug 2023)-img/p1_chem_s1_q02_B.png}{66.1}{0}
\qfigdef{p1_chem_s1_q02_C}{SRG (Enthusiast) - Test 27 (13 Aug 2023)-img/p1_chem_s1_q02_C.png}{62.5}{0}
\qfigdef{p1_chem_s1_q02_D}{SRG (Enthusiast) - Test 27 (13 Aug 2023)-img/p1_chem_s1_q02_D.png}{51.2}{0}
\qfigdef{p1_chem_s1_q05_A}{SRG (Enthusiast) - Test 27 (13 Aug 2023)-img/p1_chem_s1_q05_A.png}{32.5}{0}
\qfigdef{p1_chem_s1_q05_B}{SRG (Enthusiast) - Test 27 (13 Aug 2023)-img/p1_chem_s1_q05_B.png}{36.7}{0}
\qfigdef{p1_chem_s1_q05_C}{SRG (Enthusiast) - Test 27 (13 Aug 2023)-img/p1_chem_s1_q05_C.png}{41.8}{0}
\qfigdef{p1_chem_s1_q05_D}{SRG (Enthusiast) - Test 27 (13 Aug 2023)-img/p1_chem_s1_q05_D.png}{96.0}{0}
\qfigdef{p1_chem_s1_q07_A}{SRG (Enthusiast) - Test 27 (13 Aug 2023)-img/p1_chem_s1_q07_A.png}{137.0}{0}
\qfigdef{p1_chem_s1_q11_1}{SRG (Enthusiast) - Test 27 (13 Aug 2023)-img/p1_chem_s1_q11_1.png}{48.2}{0}
\qfigdef{p1_chem_s1_q11_A}{SRG (Enthusiast) - Test 27 (13 Aug 2023)-img/p1_chem_s1_q11_A.png}{78.9}{0}
\qfigdef{p1_chem_s1_q11_B}{SRG (Enthusiast) - Test 27 (13 Aug 2023)-img/p1_chem_s1_q11_B.png}{97.0}{0}
\qfigdef{p1_chem_s1_q14_A}{SRG (Enthusiast) - Test 27 (13 Aug 2023)-img/p1_chem_s1_q14_A.png}{99.7}{0}
\qfigdef{p1_chem_s1_q14_B}{SRG (Enthusiast) - Test 27 (13 Aug 2023)-img/p1_chem_s1_q14_B.png}{98.4}{0}
\qfigdef{p1_chem_s1_q14_C}{SRG (Enthusiast) - Test 27 (13 Aug 2023)-img/p1_chem_s1_q14_C.png}{101.7}{0}
\qfigdef{p1_chem_s1_q14_4}{SRG (Enthusiast) - Test 27 (13 Aug 2023)-img/p1_chem_s1_q14_4.png}{101.1}{0}
\qfigdef{p1_chem_s1_q17_1}{SRG (Enthusiast) - Test 27 (13 Aug 2023)-img/p1_chem_s1_q17_1.png}{161.5}{0}
\qfigdef{p1_chem_s1_q20_1}{SRG (Enthusiast) - Test 27 (13 Aug 2023)-img/p1_chem_s1_q20_1.png}{138.2}{0}
\qfigdef{p1_chem_s1_q20_A}{SRG (Enthusiast) - Test 27 (13 Aug 2023)-img/p1_chem_s1_q20_A.png}{53.9}{0}
\qfigdef{p1_chem_s1_q20_B}{SRG (Enthusiast) - Test 27 (13 Aug 2023)-img/p1_chem_s1_q20_B.png}{52.9}{0}
\qfigdef{p1_chem_s1_q20_C}{SRG (Enthusiast) - Test 27 (13 Aug 2023)-img/p1_chem_s1_q20_C.png}{45.5}{0}
\qfigdef{p1_chem_s1_q20_D}{SRG (Enthusiast) - Test 27 (13 Aug 2023)-img/p1_chem_s1_q20_D.png}{52.7}{0}
\qfigdef{p1_chem_s2_q09_A}{SRG (Enthusiast) - Test 27 (13 Aug 2023)-img/p1_chem_s2_q09_A.png}{42.3}{0}
\qfigdef{p1_chem_s2_q09_B}{SRG (Enthusiast) - Test 27 (13 Aug 2023)-img/p1_chem_s2_q09_B.png}{39.0}{0}
\qfigdef{p1_chem_s2_q09_C}{SRG (Enthusiast) - Test 27 (13 Aug 2023)-img/p1_chem_s2_q09_C.png}{43.6}{0}
\qfigdef{p1_chem_s2_q09_D}{SRG (Enthusiast) - Test 27 (13 Aug 2023)-img/p1_chem_s2_q09_D.png}{45.5}{0}
\qfigdef{p1_chem_s2_q09_5}{SRG (Enthusiast) - Test 27 (13 Aug 2023)-img/p1_chem_s2_q09_5.png}{46.8}{0}
\qfigdef{p1_chem_s2_q09_7}{SRG (Enthusiast) - Test 27 (13 Aug 2023)-img/p1_chem_s2_q09_7.png}{46.1}{0}
\qfigdef{p2_phy_s1_q03_1}{SRG (Enthusiast) - Test 27 (13 Aug 2023)-img/p2_phy_s1_q03_1.png}{76.2}{0}
\qfigdef{p2_phy_s1_q04_1}{SRG (Enthusiast) - Test 27 (13 Aug 2023)-img/p2_phy_s1_q04_1.png}{82.0}{0}
\qfigdef{p2_phy_s1_q04_2}{SRG (Enthusiast) - Test 27 (13 Aug 2023)-img/p2_phy_s1_q04_2.png}{99.2}{0}
\qfigdef{p2_phy_s2_q05_1}{SRG (Enthusiast) - Test 27 (13 Aug 2023)-img/p2_phy_s2_q05_1.png}{62.0}{0}
\qfigdef{p2_phy_s2_q06_1}{SRG (Enthusiast) - Test 27 (13 Aug 2023)-img/p2_phy_s2_q06_1.png}{86.5}{0}
\qfigdef{p2_phy_s2_q07_1}{SRG (Enthusiast) - Test 27 (13 Aug 2023)-img/p2_phy_s2_q07_1.png}{110.6}{0}
\qfigdef{p2_phy_s3_q09_1}{SRG (Enthusiast) - Test 27 (13 Aug 2023)-img/p2_phy_s3_q09_1.png}{150.1}{0}
\qfigdef{p2_phy_s3_q10_1}{SRG (Enthusiast) - Test 27 (13 Aug 2023)-img/p2_phy_s3_q10_1.png}{89.7}{0}
\qfigdef{p2_phy_s3_q11_1}{SRG (Enthusiast) - Test 27 (13 Aug 2023)-img/p2_phy_s3_q11_1.png}{79.3}{0}
\qfigdef{p2_phy_s4_q01_1}{SRG (Enthusiast) - Test 27 (13 Aug 2023)-img/p2_phy_s4_q01_1.png}{146.5}{0}
\qfigdef{p2_phy_s4_q02_1}{SRG (Enthusiast) - Test 27 (13 Aug 2023)-img/p2_phy_s4_q02_1.png}{175.6}{0}
\qfigdef{p2_phy_s4_q05_1}{SRG (Enthusiast) - Test 27 (13 Aug 2023)-img/p2_phy_s4_q05_1.png}{114.4}{0}
\qfigdef{p2_phy_s4_q06_1}{SRG (Enthusiast) - Test 27 (13 Aug 2023)-img/p2_phy_s4_q06_1.png}{159.2}{0}
\qfigdef{p2_chem_s1_q01_1}{SRG (Enthusiast) - Test 27 (13 Aug 2023)-img/p2_chem_s1_q01_1.png}{276.3}{1}
\qfigdef{p2_chem_s2_q07_1}{SRG (Enthusiast) - Test 27 (13 Aug 2023)-img/p2_chem_s2_q07_1.png}{127.3}{0}
\qfigdef{p2_chem_s3_q08_1}{SRG (Enthusiast) - Test 27 (13 Aug 2023)-img/p2_chem_s3_q08_1.png}{122.7}{0}
\qfigdef{p2_chem_s3_q08_2}{SRG (Enthusiast) - Test 27 (13 Aug 2023)-img/p2_chem_s3_q08_2.png}{151.3}{0}
\qfigdef{p2_chem_s3_q08_3}{SRG (Enthusiast) - Test 27 (13 Aug 2023)-img/p2_chem_s3_q08_3.png}{156.8}{0}
\qfigdef{p2_chem_s3_q08_4}{SRG (Enthusiast) - Test 27 (13 Aug 2023)-img/p2_chem_s3_q08_4.png}{137.7}{0}
\qfigdef{p2_chem_s3_q11_1}{SRG (Enthusiast) - Test 27 (13 Aug 2023)-img/p2_chem_s3_q11_1.png}{80.4}{0}
\qfigdef{p2_chem_s3_q11_2}{SRG (Enthusiast) - Test 27 (13 Aug 2023)-img/p2_chem_s3_q11_2.png}{72.5}{0}
\qfigdef{p2_chem_s3_q11_3}{SRG (Enthusiast) - Test 27 (13 Aug 2023)-img/p2_chem_s3_q11_3.png}{79.8}{0}
\qfigdef{p2_chem_s3_q11_4}{SRG (Enthusiast) - Test 27 (13 Aug 2023)-img/p2_chem_s3_q11_4.png}{48.7}{0}
\qfigdef{p2_chem_s3_q11_5}{SRG (Enthusiast) - Test 27 (13 Aug 2023)-img/p2_chem_s3_q11_5.png}{118.7}{0}
\qfigdef{p2_chem_s4_q03_1}{SRG (Enthusiast) - Test 27 (13 Aug 2023)-img/p2_chem_s4_q03_1.png}{129.6}{0}
\qfigdef{p2_chem_s4_q04_1}{SRG (Enthusiast) - Test 27 (13 Aug 2023)-img/p2_chem_s4_q04_1.png}{144.4}{0}
\qfigdef{p2_chem_s4_q04_2}{SRG (Enthusiast) - Test 27 (13 Aug 2023)-img/p2_chem_s4_q04_2.png}{102.6}{0}

\begin{document}
\begin{center}{\Large\bfseries SRG (Enthusiast) --- Test 27 (13 Aug 2023)}\end{center}
\begin{multicols}{2}
\papersec{1}
\partsec{1}{PHYSICS}
\sect{Section-I}
\qn{1} Calculate the equivalent capacitance between A and B.
\qfig{p1_phy_s1_q01_A}
\begin{qoptions}
\opt{A}{$\dfrac{3+\sqrt{17}}{2}C$}
\opt{B}{$\dfrac{5+\sqrt{17}}{2}C$}
\opt{C}{$\dfrac{1+\sqrt5}{2}C$}
\opt{D}{$\dfrac{3+\sqrt5}{2}C$}
\end{qoptions}
\qn{2} A point mass of 10 g oscillates according to equation $x=3\sin\left(\dfrac{\pi t}{12}+\dfrac{\pi}{4}\right)$ (here $x$ is in cm and $t$ is in s). At what time force on particle is maximum and what is the value of maximum force\,? (take $\pi^2=10$)
\begin{qoptions}
\opt{A}{15 s, $2.08\times10^{-3}$ N}
\opt{B}{18 s, $2.08\times10^{-5}$ N}
\opt{C}{15 s, $2.08\times10^{-5}$ N}
\opt{D}{18 s, $2.08\times10^{-3}$ N}
\end{qoptions}
\qn{3} \textbf{Assertion (A) :} A capacitor is connected to an ideal battery. A dielectric slab is slowly inserted into the gap between the plates of capacitor. Work required in the process is positive.\\
\textbf{Reason (R) :} Force due to capacitor will try to pull dielectric into the capacitor. To slowly insert external force needs to be outward i.e. opposite to displacement.
\begin{qoptions}
\opt{A}{\textbf{A} is not correct but \textbf{R} is correct.}
\opt{B}{Both \textbf{A} and \textbf{R} are correct but \textbf{R} is NOT the correct explanation of \textbf{A}.}
\opt{C}{Both \textbf{A} and \textbf{R} are correct and \textbf{R} is the correct explanation of \textbf{A}.}
\opt{D}{\textbf{A} is correct but \textbf{R} is not correct.}
\end{qoptions}
\qn{4} Two concentric conducting spheres are as shown. Inner sphere is a hollow shell having charge $Q$ and outer sphere is a grounded thick hollow sphere. Which of the following is correct graph of potential as a function of distance from centre?
\qfig{p1_phy_s1_q04_1}
\begin{qoptions}
\opt{A}{\qfigi{p1_phy_s1_q04_A}}
\opt{B}{\qfigi{p1_phy_s1_q04_B}}
\opt{C}{\qfigi{p1_phy_s1_q04_C}}
\opt{D}{\qfigi{p1_phy_s1_q04_D}}
\end{qoptions}
\qn{5} Which of the following is incorrect for a charge particle in a cyclotron\,?
\begin{qoptions}
\opt{A}{Particle has same frequency as of oscillator used to apply potential difference.}
\opt{B}{Particle accelerates between two dees due to electric field}
\opt{C}{Particle does not speed up in dees.}
\opt{D}{Particle accelerates only between two dees.}
\end{qoptions}
\qn{6} A particle having charge of $10^{-9}$ nC is rotating in a circle of radius $1\,\text{\AA}$ with a speed of $10^5$ m/s. What will be the magnitude of magnetic field at centre\,?
\qfig{p1_phy_s1_q06_1}
\begin{qoptions}
\opt{A}{1 T}
\opt{B}{0.1 T}
\opt{C}{10 T}
\opt{D}{such situation is not possible}
\end{qoptions}
\qn{7} What will be the magnitude of magnetic flux density at point P as shown\,?
\qfig{p1_phy_s1_q07_1}
\begin{qoptions}
\opt{A}{$\dfrac{\mu_0 i}{4\pi r}(2-\sin\theta)$}
\opt{B}{$\dfrac{\mu_0 i}{4\pi r}(2+\sin\theta)$}
\opt{C}{$\dfrac{\mu_0 i}{4\pi r}\left(\dfrac{2-\sin\theta}{\cos\theta}\right)$}
\opt{D}{$\dfrac{\mu_0 i}{4\pi r}\left(\dfrac{1-\sin\theta+\cos\theta}{\cos\theta}\right)$}
\end{qoptions}
\qn{8} Equation of motion for a point mass is $2\ddot x+5\dot x+3x=0$.\\
If all values are in S.I. unit then it is undergoing\,--
\begin{qoptions}
\opt{A}{Underdamped oscillation}
\opt{B}{Critical oscillation}
\opt{C}{Overdamped oscillation}
\opt{D}{Forced oscillation}
\end{qoptions}
\qn{9} A satellite is given velocity $v_0$ from height $h=2R$ above the earth surface as shown. Here $r$ is radius of earth. List-I has four values of speed $v_0$ and List-II has possible nature of path followed by satellite.\\
(Mass of earth $=M$)
\qfig{p1_phy_s1_q09_1}
\begin{center}\small
\begin{tabular}{|c|c||c|l|}\hline
\multicolumn{2}{|c||}{List-I} & \multicolumn{2}{c|}{List-II}\\\hline
1 & $\sqrt{\dfrac{GM}{4R}}$ & P & Circular\\\hline
2 & $\sqrt{\dfrac{GM}{3R}}$ & Q & Elliptical\\\hline
3 & $\sqrt{\dfrac{2GM}{3R}}$ & R & Parabolic\\\hline
4 & $\sqrt{\dfrac{3GM}{4R}}$ & S & Hyperbolic\\\hline
\end{tabular}
\end{center}
Which of the following is correct\,?
\begin{qoptions}
\opt{A}{$(1\to Q);\ (2\to R);\ (3\to R);\ (4\to S)$}
\opt{B}{$(1\to Q);\ (2\to P);\ (3\to R);\ (4\to S)$}
\opt{C}{$(1\to P);\ (2\to Q);\ (3\to R);\ (4\to S)$}
\opt{D}{$(1\to Q);\ (2\to Q);\ (3\to Q);\ (4\to S)$}
\end{qoptions}
\qn{10} A parallel plate capacitor has plates of dimension $10\text{ cm}\times20$ cm. Separation between plates is 4 mm. A dielectric slab of dimension $5\text{ cm}\times5\text{ cm}\times2$ mm is inserted in capacitor so that it is fully inside. If dielectric constant of slab is 9 then the capacitance will be nearly
\begin{qoptions}
\opt{A}{23 pF}
\opt{B}{49 pF}
\opt{C}{69 pF}
\opt{D}{It is not possible to calculate capacitance unless location of dielectric inside capacitor is known.}
\end{qoptions}
\qn{11} A uniform magnetic field `B' is present in a square region of side $\ell$. A point charge $q$ of mass $m$ is projected perpendicular to magnetic field with a speed of $\dfrac{qB\ell}{2m}$ from a point on boundary. What is the maximum duration for which particle can be inside field\,?
\qfig{p1_phy_s1_q11_1}
\begin{qoptions}
\opt{A}{$\dfrac{\pi m}{2qB}$}
\opt{B}{$\dfrac{\pi m}{qB}$}
\opt{C}{$\dfrac{3\pi m}{2qB}$}
\opt{D}{$\dfrac{2\pi m}{qB}$}
\end{qoptions}
\qn{12} A planet has density $4\times10^3$ kg/m$^3$. A tunnel is dug along the diameter of the planet. How much time will it take for a particle to move from one end to another end of tunnel under the action of only gravitational force?
\begin{qoptions}
\opt{A}{42.3 min}
\opt{B}{49.5 min}
\opt{C}{84.6 min}
\opt{D}{57.2 min}
\end{qoptions}
\qn{13} Two identical particles each of which has mass $m$ and charge $q$ are kept at some separation. What is the ratio of $\dfrac qm$ if gravitational potential energy is equal to electrostatic potential energy\,?
\begin{qoptions}
\opt{A}{$8.6\times10^{-11}$ C/kg}
\opt{B}{$7.4\times10^{-25}$ C/kg}
\opt{C}{$7.4\times10^{-11}$ C/kg}
\opt{D}{It is not possible that energies are equal.}
\end{qoptions}
\qn{14} Work required to bring a $-2\,\mu$C charge particle from $\infty$ to P is 2 mJ then what will be work done by electric field to take $3\,\mu$C charge from P to infinity\,?
\qfig{p1_phy_s1_q14_1}
\begin{qoptions}
\opt{A}{3 mJ}
\opt{B}{$-3$ mJ}
\opt{C}{1 mJ}
\opt{D}{$-1$ mJ}
\end{qoptions}
\qn{15} Two concentric conducting shells can be connected by a switch `s' as shown. Inner shell has charge $Q$.\\
\textbf{Statement-1 :} Charge on innear sphere flows from inner shell to outer shell when switch is closed.\\
\textbf{Statement-2 :} Charge will not flow it $Q$ is negative as charge can not flow from low to high potential.
\qfig{p1_phy_s1_q15_1}
\begin{qoptions}
\opt{A}{Statement-I is incorrect and statement-II is correct}
\opt{B}{Statement-I is correct and Statement-II is incorrect}
\opt{C}{Both Statement-I and Statement-II are false}
\opt{D}{Both Statement-I and Statement-II are true}
\end{qoptions}
\qn{16} A square frame of side `$\ell$' is made of four rods. Linear charge density of rods are shown in figure. Direction of electric field at centre is along--
\qfig{p1_phy_s1_q16_1}
\begin{qoptions}
\opt{A}{CP}
\opt{B}{CQ}
\opt{C}{CR}
\opt{D}{CS}
\end{qoptions}
\qn{17} At a certain point on earth declination is $+30^\circ$ and dip angle is $+60^\circ$. Horizontal component of magnetic field at that point is 0.12G. What will be net magnetic field in geographical meridian\,?
\begin{qoptions}
\opt{A}{0.10 G}
\opt{B}{0.20 G}
\opt{C}{0.23 G}
\opt{D}{None of these}
\end{qoptions}
\qn{18} Consider the following four statements.\\
1.\ \ A paramagnetic substance displays greater magnetisation when cooled.\\
2.\ Diamagnetism is independent of temperature.\\
3.\ If a solenoid sues copper for its core instead of air then field will be greater.\\
4.\ Magnetic field lines are always perpendicular to surface of ferromagnetic material.\\
Which statements are true?
\begin{qoptions}
\opt{A}{Only 1 and 2}
\opt{B}{1, 2 and 3}
\opt{C}{1, 2, and 4}
\opt{D}{All four}
\end{qoptions}
\qn{19} A car battery supplies a high current to starter motor for a short duration. Current in the wire is 250 A and separation between wires is 0.5 cm. What will be the force per unit length between wires\,?
\begin{qoptions}
\opt{A}{1.25 N/m}
\opt{B}{1.625 N/m}
\opt{C}{2.5 N/m}
\opt{D}{5 N/m}
\end{qoptions}
\qn{20} Escape speed from a planet is 10 km/s. A body is projected from the surface of planet with twice of escape speed. What will be the speed of body at large distance from planet\,?
\begin{qoptions}
\opt{A}{14.14 km/s}
\opt{B}{17.32 km/s}
\opt{C}{20 km/s}
\opt{D}{30 km/s}
\end{qoptions}
\sect{Section-II}
\qn{1} Given figure shows two pendulums. Penduulum-1 is a physical pendulum made by uniform rod and pendulum-2 is simple pendulum. Both have equal mass. Length of two pendulum is $\ell$ and $\dfrac{\ell}{4}$ respectively. If both pendulums are oscillating having same amplitude then the ratio of energy of oscillation of pendulum-2 and pendulum-1 is $x$. Calculate $64x$.
\qfig{p1_phy_s2_q01_1}
\qn{2} An infinitely long wire carrying current $I$ is kept perpendicular to the plane of paper as shown. $B$ is magnetic field due to wire. Current in the wire is in outward direction. If $\displaystyle\int_P^Q\vec B\cdot d\vec\ell=\dfrac{\mu_0 I}{x}$ then $x=$ \underline{\hspace{1.5cm}}.
\qfig{p1_phy_s2_q02_1}
\qn{3} A ballistic galvanometer, which measures charge flow through it, is connected in the circuit as shown. All capacitors are initially uncharged. First switch `s$_1$' is closed and reading is noted. Then witch `s$_2$' then is closed and again reading is noted. The difference in reading is \underline{\hspace{1cm}} $\mu$C. ($V=36$ Volt, $C=4\,\mu$F)
\qfig{p1_phy_s2_q03_1}
\qn{4} 6 stars of equal mass are orbiting in circular path around their centre of mass `c' as shown. If the speed of any one of the star is $\sqrt{\dfrac{GM}{r}\left(\dfrac{\alpha}{4}+\dfrac{1}{\sqrt\beta}\right)}$, then $\alpha+\beta$ is
\qfig{p1_phy_s2_q04_1}
\qn{5} Charge density on upper plate of capacitor 18 nC/m$^2$. What will be the bounded charge density (in nC/m$^2$) on upper surface of dielectric of dielectric constant 3\,?
\qfig{p1_phy_s2_q05_1}
\qn{6} A conducting sphere of radius $R$ is grounded as shown. A ring of radius $r$ is charged such that potential at its centre is 90 V. Ring is now placed such that the centre of ring is at the top of sphere. Find the total charge induced (in pC) on sphere. [$R=4$ cm and $4r=3R$]
\qfig{p1_phy_s2_q06_1}
\qn{7} Two identical capacitors of $2\mu$F capacitance are connected in series as shown. One of the capacitor has initial potential difference across it as shown. Now switch is closed. When circuit reaches steady state a dielectric of dielectric constant 3 inserted into right capacitor. If final potential on left capacitor is $x$ volt, then $16x=$ \underline{\hspace{1.5cm}}
\qfig{p1_phy_s2_q07_1}
\qn{8} Given circuit is in steady state. Now switch `S' is closed. Charge flowing through the switch is \underline{\hspace{1.5cm}} $\mu$C.
\qfig{p1_phy_s2_q08_1}
\qn{9} In the figure shown a point charge $2Q$ is placed at the centre. The two concentric conducting spherical shells of radii `$R$' and `$2R$' shown have charges $+Q$ and $-2Q$ respectively. When both shells are connected together by a connecting wire, $xQ$ charge flows to outer shell. Calculate $x$.
\qfig{p1_phy_s2_q09_1}
\qn{10} A solid conducting sphere of radius $R=20$ cm has two spherical cavities of radius $a=3$ cm and $b=4$ cm. The cavities have charges $q_a$ \& $q_b$ respectively at their centres. Centers of cavities are at 10 cm away from center of sphere as shown. What will be the magnitude of force (in N) on $q_a$ due to $q_b$\,?
\qfig{p1_phy_s2_q10_1}
\partsec{2}{CHEMISTRY}
\sect{Section-I}
\qn{1} \textbf{INCORRECT} order of stability of complexes
\begin{qoptions}
\opt{A}{$[\mathrm{Cu(trien)}]^{++}>[\mathrm{Cu(en)_2}]^{++}$}
\opt{B}{$[\mathrm{Ni(ox)_3}]<[\mathrm{Ni(acac)_3}\,]^+$}
\opt{C}{$[\mathrm{Ni(H_2O)_6}]^{++}>[\mathrm{NiF_6}]^{4-}$}
\opt{D}{$[\mathrm{CoF_3(H_2O)_3}]>[\mathrm{CoF_6}]^{3-}$}
\end{qoptions}
\qn{2} For which of the following enol content is more than the Keto.
\begin{qoptions}
\opt{A}{\qfigi{p1_chem_s1_q02_A}}
\opt{B}{\qfigi{p1_chem_s1_q02_B}}
\opt{C}{\qfigi{p1_chem_s1_q02_C}}
\opt{D}{\qfigi{p1_chem_s1_q02_D}}
\end{qoptions}
\qn{3} \textbf{Assertion:} The entropy change as well as the gibb's free energy change both have non-zero value for the formation of an ideal solution by mixing two different substances.\\
\textbf{Reason :} Enthalpy remains constant during the formation of an ideal solution by mixing two different substances.
\begin{qoptions}
\opt{A}{I \& II both are correct but II is not a correct explanation of I.}
\opt{B}{I \& II both are correct but II is a correct explanation of I.}
\opt{C}{I is correct but II is incorrect.}
\opt{D}{I is incorrect but II is incorrect.}
\end{qoptions}
\qn{4} The correct order for the mentioned property is:
\begin{qoptions}
\opt{A}{Be$^{2+}<$ Al$^{3+}$ (Size)}
\opt{B}{Cd $<$ In (Ionisation Energy)}
\opt{C}{S $<$ Se (Electron Affinity)}
\opt{D}{O $<$ Cl (Electronegativity)}
\end{qoptions}
\qn{5} Rate of S$_\mathrm{N}$1 is the maximum in
\begin{qoptions}
\opt{A}{\qfigi{p1_chem_s1_q05_A}}
\opt{B}{\qfigi{p1_chem_s1_q05_B}}
\opt{C}{\qfigi{p1_chem_s1_q05_C}}
\opt{D}{\qfigi{p1_chem_s1_q05_D}}
\end{qoptions}
\qn{6} A certain solid that is soft, nonmalleable, good conductor of electricity and has high melting point can belong to which of the following category of solids
\begin{qoptions}
\opt{A}{Molecular solid}
\opt{B}{Metallic Solid}
\opt{C}{Covalent solid}
\opt{D}{Ionic solid}
\end{qoptions}
\qn{7} Consider:
\qfig{p1_chem_s1_q07_A}
\begin{qoptions}
\opt{A}{A \& B are enantiomers}
\opt{B}{A \& B are diastereomers}
\opt{C}{A \& B are Identical}
\opt{D}{A \& B are conformers}
\end{qoptions}
\qn{8} Which pair of molecules have the same shape but different hybridization of the central atom
\begin{qoptions}
\opt{A}{XeF$_6$, IF$_7$}
\opt{B}{I$_3^-$, I$_3^+$}
\opt{C}{CO$_2$, XeF$_2$}
\opt{D}{CF$_4$, SF$_4$}
\end{qoptions}
\qn{9} Match the column
\begin{center}\small
\begin{tabular}{|p{0.4\linewidth}|p{0.45\linewidth}|}\hline
Column (I) (Aqueous solutions of density 1.2 gml$^{-1}$) & Column (II) Facts about the weights of the components of the solutions\\\hline
(i) 1L of 5M NaOH(aq) & (P) W$_\text{solvent}=1$ kg\\\hline
(ii) 1 kg of 5m NaOH(aq) & (Q) W$_\text{solvent}>1$ kg\\\hline
(iii) 1kg of NaOH(aq) $X_\text{NaOH}=0.5$ & (R) W$_\text{solute}>200$g\\\hline
(iv) 1L of NaOH Strength=5 g/l & (S) W$_\text{solute}<200$g\\\hline
\end{tabular}
\end{center}
\begin{qoptions}
\opt{A}{i $\to$ P, ii$\to$ S, iii $\to$ Q, iv $\to$ R}
\opt{B}{i $\to$ Q, ii$\to$ P, iii $\to$ S, iv $\to$ R}
\opt{C}{i $\to$ Q, ii$\to$ P, iii $\to$ R, iv $\to$ S}
\opt{D}{i $\to$ P, ii$\to$ S, iii $\to$ R, iv $\to$ Q}
\end{qoptions}
\qn{10} The IUPAC name of that hydrate isomer of CrCl$_3$.6H$_2$O , which has the lowest possible electrical conductivity in its aqueous solution.
\begin{qoptions}
\opt{A}{Hexaaquachromium(III) chloride}
\opt{B}{Tetraaquadichloridochromium(II) chloridedihydrate}
\opt{C}{Pentaaquachloridochromium(III) chloridemonohydrate}
\opt{D}{Triaquatrichloridochromium (III) trihydratc}
\end{qoptions}
\qn{11} \qfigi{p1_chem_s1_q11_1} can be prepared by:
\begin{qoptions}
\opt{A}{\qfigi{p1_chem_s1_q11_A}}
\opt{B}{\qfigi{p1_chem_s1_q11_B}}
\opt{C}{$\mathrm{CH_3-CH=O}\ \xrightarrow[(2)\,\mathrm{H^+}]{(1)\,\mathrm{CH_3MgBr}}$}
\opt{D}{$\mathrm{CH_3-CO-OC_2H_5}\ \xrightarrow[(2)\,\mathrm{H^+}]{(1)\,\mathrm{LiAlH_4}}$}
\end{qoptions}
\qn{12} Which defect is not associated with the frenkel defects in ionic solids
\begin{qoptions}
\opt{A}{Vacancy defect}
\opt{B}{Interstitial defect}
\opt{C}{Dislocation effect}
\opt{D}{Substitution defect}
\end{qoptions}
\qn{13} Which of the following has all the bonds non-polar but still the molecule is polar
\begin{qoptions}
\opt{A}{N$_2$}
\opt{B}{O$_3$}
\opt{C}{P$_4$}
\opt{D}{None of these}
\end{qoptions}
\qn{14} Saytzeff's alkene is the major product in
\begin{qoptions}
\opt{A}{\qfigi{p1_chem_s1_q14_A}}
\opt{B}{\qfigi{p1_chem_s1_q14_B}}
\opt{C}{\qfigi{p1_chem_s1_q14_C}}
\opt{D}{\qfigi{p1_chem_s1_q14_4}}
\end{qoptions}
\qn{15} \textbf{Statement 1:} The osmotic pressure of a solution is the excess pressure that must be applied to a solution to prevent osmosis, i.e., to stop the net passage of solvent molecules through a semipermeable membrane into the solution.\\
\textbf{Statement. 2:} During osmosis the solvent particles move from the region of the higher solvent concentration to the region of lower solvent concentration, across a semipermeable membrane
\begin{qoptions}
\opt{A}{Statement I is true but Statement II is false}
\opt{B}{Statement I is false but Statement II is true}
\opt{C}{Both Statement I and Statement II are true}
\opt{D}{Both Statement I and Statement II are false}
\end{qoptions}
\qn{16} Which of the following has magnesium in it?
\begin{qoptions}
\opt{A}{Chlorophyll}
\opt{B}{Hemoglobin}
\opt{C}{Carbonic anhydrase}
\opt{D}{Vitamin B$_{12}$}
\end{qoptions}
\qn{17} Consider the following reactions where, $k_1$, $k_2$, $k_3$ and $k_4$ are the rate constants for the reactions (1), (2), (3), (4) respectively.
\qfig{p1_chem_s1_q17_1}
Choose the CORRECT options regarding the rate constants:
\begin{qoptions}
\opt{A}{$k_1>k_2,\ k_3>k_4$}
\opt{B}{$k_1<k_2,\ k_3>k_4$}
\opt{C}{$k_1>k_2,\ k_3<k_4$}
\opt{D}{$k_1<k_2,\ k_3>k_4$}
\end{qoptions}
\qn{18} Consider the statements:\\
(P) All the solutions that show any positive deviation from the ideal behaviour, must form a maximum boiling azeotrope\\
(Q) All the solutions that show any positive deviation from the ideal behaviour, must form a minimum boiling azeotrope\\
(R) Solution of nitric acid and water is an example of maximum boiling azeotrope. This azeotrope has the approximate composition, 68\% nitric acid and 32\% water by mass, with a boiling point of 393.5 K.\\
The correct order of true (T) or false (F.) statements is.
\begin{qoptions}
\opt{A}{FFF}
\opt{B}{TFT}
\opt{C}{FTT}
\opt{D}{FFT}
\end{qoptions}
\qn{19} Which of the following pairs of orbitals (atomic as well as molecular) have the same number of nodal plane(s)
\begin{qoptions}
\opt{A}{$2p_z$ and $(2p_z-2p_z)\,\sigma$}
\opt{B}{$(2p_z-2p_z)\sigma^*$ and $(3d_{xy}-3d_{xy})\,\pi$}
\opt{C}{$(1s-1s)\sigma^*$ and $(3d_{xy}-3d_{xy})\,\pi^*$}
\opt{D}{$1s$ and $3d_{xy}$}
\end{qoptions}
\qn{20} \qfig{p1_chem_s1_q20_1}
product P is\,:
\begin{qoptions}
\opt{A}{\qfigi{p1_chem_s1_q20_A}}
\opt{B}{\qfigi{p1_chem_s1_q20_B}}
\opt{C}{\qfigi{p1_chem_s1_q20_C}}
\opt{D}{\qfigi{p1_chem_s1_q20_D}}
\end{qoptions}
\sect{Section-II}
\qn{1} An organic compound has 48\% each of carbon and oxygen, the rest being hydrogen. Also, One mole of the compound has just sufficient oxygen needed for the complete combustion of 48 gram of methane. Find the gram molecular mass of the organic compound in gram?
\qn{2} For a certain wavelengths, $\lambda$, at a temperature t$^\circ$C, the specific angle of rotation of the plane polarized right is $+20^\circ$ for a sample of pure d-2-chlorobutane. What is the percentage of the `$\ell$' isomer in a sample of 2-chlorobutane exhibiting the specific angle of rotation to be `$-12^\circ$' for the same $\lambda$ and t?
\qn{3} The maximum number of electrons in the outermost shell of any element ( from among all the 118 known elements ) is equal to :
\qn{4} How many of the following have a higher ebullioscopic constant (molal elevation constant) than water:\\
Ethanol, Cyclohexane, Benzene, Chloroform, Carbon tetrachloride, Carbon disulphide, Diethyl ether, Acetic acid
\qn{5} How many of the following molecules are planar\\
ICl$_3$, \quad I$_2$Cl$_6$, \quad AlCl$_3$, \quad Al$_2$Cl$_6$, \quad H$_2$C$=$C$=$C$=$CH$_2$
\qn{6} The mass of ice in grams that separates from a 1 kg solution of 0.1 mole urea in water when the solution is kept at $-2^\circ$C is?\\
Use, K$_\mathrm{f}=2$K mol$^{-1}$kg. Assume the solution to follow the Raoult's law all the time.
\qn{7} Let,\\
\textbf{X} = Number of Bravais crystal systems that have atleast 2 different types of lattice variations out of the possible 4 variations (Primitive, Body-Centered, Face Centered \& End-Centered)\\
\textbf{Y} = Number of Bravais crystal systems that have at least two axial angles to be the same.\\
Find $10\,X+Y$?
\qn{8} Find the value of X + Y for complex $[\mathrm{CoBr_2(CN)(NO_2)(H_2O)(NH_3)}]$\\
where: \quad X = Number of compounds when both Br$^-$ are at cis positions.\\
\hspace*{2.6em} Y = Number of compounds when both Br$^-$ are at trans positions.
\qn{9} Consider the following reactions:
\begin{center}\small
\begin{tabular}{|c|c|l|}\hline
(A) & \qfigi{p1_chem_s2_q09_A} & $\xrightarrow[(2)\,\mathrm{H^+}]{(1)\,\mathrm{CH_3MgBr}}$\\\hline
(B) & \qfigi{p1_chem_s2_q09_B} & $\xrightarrow[(2)\,\mathrm{H_2O}]{(1)\,\mathrm{HBr}}$\\\hline
(C) & \qfigi{p1_chem_s2_q09_C} & $\xrightarrow[(2)\,\mathrm{H_2O}]{(1)\,\mathrm{HBr}}$\\\hline
(D) & \qfigi{p1_chem_s2_q09_D} & $\xrightarrow[(2)\,\mathrm{H_2O}]{(1)\,\mathrm{HBr}}$\\\hline
(E) & \qfigi{p1_chem_s2_q09_5} & $\xrightarrow[(2)\,\mathrm{CH_2O},\ (3)\,\mathrm{H_3O^+}]{(1)\,\mathrm{Mg/Ether}}$\\\hline
(F) & \qfigi{p1_chem_s2_q09_7} & $\xrightarrow[(2)\,\mathrm{H_2O}]{(1)\,\mathrm{HBr/Peroxide}}$\\\hline
\end{tabular}
\end{center}
How many can produce 1-methylcyclohexanol?
\qn{10} Find the Effective atomic number of $[\mathrm{Fe(CO)_2}\,(\pi\text{--}\mathrm{C_5H_5})\mathrm{CI}]$.
\partsec{3}{MATHEMATICS}
\sect{Section-I}
\qn{1} If $f(x)=kx^2-m$ such that $f(1)\in[-4,-1]$ and $f(2)\in[-1,5]$ then the max. value of $f(3)$ is
\begin{qoptions}
\opt{A}{15}
\opt{B}{19}
\opt{C}{20}
\opt{D}{25}
\end{qoptions}
\qn{2} Let `f' be a function such that\\
$4f\left(\dfrac{1+x}{x}\right)+8f(x+1)=\log_{12}x$ then the value of $4\bigl(f(10)+f(13)+f(17)\bigr)$
\begin{qoptions}
\opt{A}{5}
\opt{B}{2}
\opt{C}{10}
\opt{D}{3}
\end{qoptions}
\qn{3} \textbf{Assertion:} $\sin^{-1}\left(\dfrac{1}{\sqrt e}\right)>\tan^{-1}\left(\dfrac{1}{\sqrt\pi}\right)$\\
\textbf{Reason:} $\sin^{-1}x>\tan^{-1}y$ for $x>y,\ \forall\ x,y\in(0,1)$
\begin{qoptions}
\opt{A}{Both assertion and reason are true and reason is correct explanation of assertion}
\opt{B}{Both assertion and reason are true and reason is not the correct explanation of assertion}
\opt{C}{Assertion is true and reason is false}
\opt{D}{Both assertion and reason are false}
\end{qoptions}
\qn{4} The largest integral value of $x$ for which $\cos^{-1}(\cos4)>3x^2-4x$ is
\begin{qoptions}
\opt{A}{0}
\opt{B}{1}
\opt{C}{2}
\opt{D}{3}
\end{qoptions}
\qn{5} Number of common point for the curves\\
$y=\sin^{-1}(2x)+\tan^{-1}\left(\dfrac{1}{[2x]}\right)+2$ and\\
$y=\cos^{-1}(2x+5)+1$ is (where $[.]$ denotes greatest integer function)
\begin{qoptions}
\opt{A}{0}
\opt{B}{1}
\opt{C}{3}
\opt{D}{4}
\end{qoptions}
\qn{6} Let $f(x)=\dfrac{\log\left(1+x+x^2\right)+\log\left(1-x+x^2\right)}{\sec x-\cos x}$, $x\ne0$. Then the value of $f(0)$ so that $f$ is continuous at $x=0$ is
\begin{qoptions}
\opt{A}{1}
\opt{B}{0}
\opt{C}{3}
\opt{D}{$-1$}
\end{qoptions}
\qn{7} If $a$, $b$, $c$ and $d$ are real numbers such that $\dfrac{3a+2b}{c+d}+\dfrac32=0$ then the equation $ax^3+bx^2+cx+d=0$ has
\begin{qoptions}
\opt{A}{at least one root in $[-2,0]$}
\opt{B}{at least one root in $[0,2]$}
\opt{C}{at least two roots in $[-2\ 2]$}
\opt{D}{No root in $[-2,2]$}
\end{qoptions}
\qn{8} If the acute angles between the curves $y=|x^2-1|$ and $y=|x^2-3|$ at their points of intersection be $\theta$ such that $\sin\theta=\dfrac{4\sqrt2}{k}$, then $k$ is equal to
\begin{qoptions}
\opt{A}{5}
\opt{B}{7}
\opt{C}{3}
\opt{D}{9}
\end{qoptions}
\qn{9} Match the following lists :
\begin{center}\small
\begin{tabular}{|p{0.62\linewidth}|c|}\hline
\multicolumn{1}{|c|}{List-I} & List-II\\\hline
1.\ Number of roots of the equation $\sin^{-1}(\sin x)=\dfrac x4$ & P.\ 2\\\hline
2.\ Number of roots of the equation $\cos^{-1}(\cos x)=\dfrac x4$ & Q.\ 3\\\hline
3.\ Number of roots of the equation $\tan^{-1}(\tan x)=\dfrac x4$ & R.\ 4\\\hline
4.\ Number of roots of the equation $\cot^{-1}(\cot x)=\dfrac x4$ & S.\ 5\\\hline
\end{tabular}
\end{center}
\begin{qoptions}
\opt{A}{$(1\to Q);\ (2\to R);\ (3\to Q);\ (4\to P)$}
\opt{B}{$(1\to Q);\ (2\to P);\ (3\to Q);\ (4\to P)$}
\opt{C}{$(1\to Q);\ (2\to R);\ (3\to S);\ (4\to P)$}
\opt{D}{$(1\to S);\ (2\to R);\ (3\to Q);\ (4\to P)$}
\end{qoptions}
\qn{10} Let $f(x)$ be differentiable for real $x$ such that\\
$f'(x)>0$ on $(-\infty,-4)$,\\
$f'(x)<0$ on $(-4,6)$ and $f'(x)>0$ on $(6,\infty)$.\\
If $g(x)=f(10-2x)$, then the value of $g'(2)$ is
\begin{qoptions}
\opt{A}{1}
\opt{B}{2}
\opt{C}{0}
\opt{D}{4}
\end{qoptions}
\qn{11} The number of functions $f$ from the set $A=\{0,1,2\}$ into the set $B=\{0,1,2,3,4,5,6,7\}$ such that $f(i)\le f(j)$ for $i<j$ and $\forall\ i,j\in A$ is
\begin{qoptions}
\opt{A}{${}^8C_3$}
\opt{B}{${}^8C_3+2({}^8C_2)$}
\opt{C}{${}^8C_3+{}^{10}C_3$}
\opt{D}{${}^{10}C_3$}
\end{qoptions}
\qn{12} Consider $f(x)=\begin{cases}2-|x| & -1\le x\le1\\ |x-2|-x & 1<x\le3\end{cases}$ and\\
$g(x)=\begin{cases}\sin x-1 & 0\le x<\frac\pi2\\ [x]-\cos(x-2) & \frac\pi2\le x\le\pi\end{cases}$\\
(where $[.]$ represents greatest integer function) then which of following is correct.
\begin{qoptions}
\opt{A}{$\displaystyle\lim_{x\to1^+}g\bigl(f(x)\bigr)=-1$}
\opt{B}{$\displaystyle\lim_{x\to0^-}g\bigl(f(x)\bigr)=0$}
\opt{C}{$\displaystyle\lim_{x\to2^+}\dfrac{f\bigl(g(x)\bigr)}{f(x)-2}=\dfrac12$}
\opt{D}{$\displaystyle\lim_{x\to0^+}\dfrac{g\bigl(f(x)\bigr)}{\bigl(f(x)-2\bigr)^2}=\dfrac{-1}{2}$}
\end{qoptions}
\qn{13} Let $f(x)=\sin x+\left(x^3-3x^2+4x-2\right)\cos x$, $\forall\ x\in(0,1)$ consider the following statement\\
Statement : I - `f' has a zero in $(0,1)$\\
Statement : II - `f' is monotonic in $(0,1)$\\
then
\begin{qoptions}
\opt{A}{I and II are true}
\opt{B}{I is true and II is false}
\opt{C}{I is false and II is true}
\opt{D}{I and II are false}
\end{qoptions}
\qn{14} The max value of $f(x)=e^x+x\ln x$ on the interval $x\in[1,2]$ is
\begin{qoptions}
\opt{A}{$e^2+\ln2+1$}
\opt{B}{$e^2+2\ln2$}
\opt{C}{$e^{\pi/2}+\dfrac\pi2\ln2$}
\opt{D}{$e^{3/2}+\dfrac32\ln\dfrac32$}
\end{qoptions}
\qn{15} Let $a,b\in\mathbb R$ such that the function $f$ given by $f(x)=\ln|x|+bx^2+ax$, $x\ne0$ has extreme values at $x=-1$ and $x=2$.\\
\textbf{Statement-1 :} $f$ has local maximum at $x=-1$ and at $x=2$.\\
\textbf{Statement-2 :} $a=\dfrac12$ and $b=\dfrac{-1}{4}$.
\begin{qoptions}
\opt{A}{Statement-1 is false, Statement-2 is true.}
\opt{B}{Statement-1 is true, statement-2 is true; statement-2 is a correct explanation for Statement-1.}
\opt{C}{Statement-1 is true, statement-2 is true; statement-2 is not a correct explanation for Statement-1.}
\opt{D}{Statement-1 is true, statement-2 is false.}
\end{qoptions}
\qn{16} Find largest value of the non-negative integer `k' if $\displaystyle\lim_{x\to4}\left(\dfrac{-kx+\sin(x-4)+4k}{x+\sin(x-4)-4}\right)^{\frac{4-x}{4-2\sqrt x}}=\dfrac14$
\begin{qoptions}
\opt{A}{0}
\opt{B}{2}
\opt{C}{1}
\opt{D}{3}
\end{qoptions}
\qn{17} Let $f(x)=e^{g(x)}$ where $f(x)$ is a twice differentiable positive function on $(0,\infty)$ i.e. $e^{f(x+1)}=\left(e^{f(x)}\right)^x$ then the value of $g''\left(\dfrac{75}{2}\right)-g''\left(\dfrac12\right)$
\begin{qoptions}
\opt{A}{$-4\left\{1+\dfrac19+\dfrac1{25}+\dots+\dfrac1{(75)^2}\right\}$}
\opt{B}{$-4\left\{1+\dfrac19+\dfrac1{25}+\dots+\dfrac1{(73)^2}\right\}$}
\opt{C}{$4\left\{1+\dfrac19+\dfrac1{25}+\dots+\dfrac1{(75)^2}\right\}$}
\opt{D}{$4\left\{1+\dfrac19+\dfrac1{25}+\dots+\dfrac1{(73)^2}\right\}$}
\end{qoptions}
\qn{18} Let $f(xy)=xf(y)+yf(x)$, $\forall\ x,y\in\mathbb R^+$ and $f(x)$ is differentiable in $(0,\infty)$, $f'(1)=2$, then the value of $f(e)$
\begin{qoptions}
\opt{A}{$e$}
\opt{B}{$2e$}
\opt{C}{$3e$}
\opt{D}{$4e$}
\end{qoptions}
\qn{19} If $f(x)$ be a continuous function on $\mathbb R$ such that $f(x)=f(8-x)$ and $f(x+10)=f(14-x)$, $\forall\ x\in\mathbb R$ and $f(0)=f(3)=6$ then minimum number of `$\alpha$' such that $f(\alpha)=6$, where $\alpha\in[0,100]$
\begin{qoptions}
\opt{A}{26}
\opt{B}{24}
\opt{C}{16}
\opt{D}{20}
\end{qoptions}
\qn{20} If the tangent at $(1,1)$ on $y^2=x(2-x)^2$ meets the curve again at P, then find coordinates of P is:
\begin{qoptions}
\opt{A}{$\left(\dfrac94,\dfrac38\right)$}
\opt{B}{$\left(\dfrac94,\dfrac32\right)$}
\opt{C}{$\left(\dfrac67,\dfrac32\right)$}
\opt{D}{$\left(\dfrac67,\dfrac{-3}{2}\right)$}
\end{qoptions}
\sect{Section-II}
\qn{1} If $x$ is real, then the maximum value of $y=2(3-x)\left(x+\sqrt{x^2+4}\right)$ is :
\qn{2} The value of {\scriptsize$\left[\displaystyle\lim_{n\to\infty}\left(5\times3^2\times5^3\times3^4\dots5^{n-1}\times3^n\right)^{\frac{1}{(n^2+1)}}\right]^4$} is:
\qn{3} If $a$ and $b$ are the number of points of non differentiability of $f(x)=[\sin^{-1}x]$, $x\in[-1,1]$ and $g(x)=\left[\dfrac{2}{1+x^2}\right]$, $x\ge0$, (where $[.]$ represents greatest integer function) respectively, then the value of $ab$ is
\qn{4} Let $f(x)=[\tan x\,[\cot x]]$, $x\in\left[\dfrac\pi{12},\dfrac\pi2\right)$ (where $[.]$ denotes the greatest integer less than or equal to $x$). Then the number of points, where $f(x)$ is discontinuous is
\qn{5} The differential coefficient of\\ {\small$\sin^{-1}\left(\dfrac{4\sin2x+3\cos2x}{5}\right)$} w.r.t.\\ {\small$\cos^{-1}\left(\dfrac{5\cos x-4\sin x}{\sqrt{41}}\right)$}\\ is $k$ then $2k$ equals :
\qn{6} If fractional part of $\dfrac1x$ and $x^2$ for some $x\in(\sqrt2,\sqrt3)$ are equal then the value of $x^4-\dfrac3x$ is :
\qn{7} Let $f:[2,15]\to\{1,2,3\}$, and the number of function $f$ can be formed on its domain such that it is discontinuous only on all integral points are $a.2^b$ (where $a,b\in N$, $a$ is prime number) then $\left[\dfrac ba\right]$ is :\\
(where $[.]$ denotes greatest integer function)
\qn{8} If $\displaystyle\lim_{x\to1}\dfrac{x^{1006}-(1006)x+1005}{xe^x-e^x-ex+e}=\ell$\\
then sum of digit of the number $\ell e$
\qn{9} If $f(x)=\left(x^3-1\right)|x-1|+\sin(\pi x)\,|x-2|+\bigl|\,||x|-3|-5\bigr|$ and $\lambda$ be the number of points in which $f(x)$ is non-differentiable then last two digit of the number $\lambda^{32}$ is :
\qn{10} If $a$ and $b$ be the least and greatest value of the function $f(x)=3\sin^{-1}x-4\cot^{-1}x$ then $\left[\dfrac{b-a}{\pi}\right]$ is :\\
(where $[.]$ denotes greatest integer function)
\papersec{2}
\partsec{1}{PHYSICS}
\sect{Section-I (i)}
\qn{1} A 1 kg mass stretches a spring by 49 cm. The system is immersed in a medium that impacts a damping force equal to four times the instantaneous velocity of the mass. Particle is released from a point 24 cm above the equilibrium. If upward direction is taken positive then equation of motion is -
\begin{qoptions}
\opt{A}{$e^{-4t}\left[0.24\cos4t+0.24\sin4t\right]$}
\opt{B}{$e^{-2t}\left[0.24\cos4t+0.24\sin4t\right]$}
\opt{C}{$e^{-2t}\left[0.24\cos4t+0.12\sin4t\right]$}
\opt{D}{$e^{-4t}\left[0.24\cos4t-0.24\sin4t\right]$}
\end{qoptions}
\qn{2} A metal sphere of inner radius $R$ and outer radius $2R$ is filled with a space charge of uniform charge density $\rho$. How much energy is released if sphere is grounded\,?
\begin{qoptions}
\opt{A}{$\dfrac{\rho^2\pi R^5}{45\varepsilon_0}$}
\opt{B}{$\dfrac{2\rho^2\pi R^5}{45\varepsilon_0}$}
\opt{C}{$\dfrac{\rho^2\pi R^5}{9\varepsilon_0}$}
\opt{D}{$\dfrac{\rho^2\pi R^5}{18\varepsilon_0}$}
\end{qoptions}
\qn{3} Two identical capacitor of capacitance 10 $\mu$F are charged to a potential 10 V as shown. At $t=0$, switch $S_1$ is closed and at $t=\ln4$ ms switch $S_2$ is also closed. Current through resistor at $t=\ln8$ ms is nearly
\qfig{p2_phy_s1_q03_1}
\begin{qoptions}
\opt{A}{44 mA}
\opt{B}{88 mA}
\opt{C}{53 mA}
\opt{D}{27 mA}
\end{qoptions}
\qn{4} A point charge is kept at the centre of two layers of spherical dielectrics as shown. Electric field strength varies as a function of distance from point charge as shown.
\qfig{p2_phy_s1_q04_1}
\qfig{p2_phy_s1_q04_2}
The dielectric constant of outer sphere is $K_2=$ \_\_\_\_\_\_\_ ($K_1=15$)
\begin{qoptions}
\opt{A}{25}
\opt{B}{2.5}
\opt{C}{12.5}
\opt{D}{1.25}
\end{qoptions}
\sect{Section-I (ii)}
\qn{5} A conical pendulum undergoes circular motion with angular speed $\dfrac52\hat{j}$ rad/s. Take point of suspension as origin. Mass of string is negligible but it has mass $m=1$ kg and charge $q=20$ mC. At $t=0$, particle is in $xy$ plane as shown, which of the following option(s) is are correct\,? [Take $g=10\ \text{ms}^{-2}$]
\qfig{p2_phy_s2_q05_1}
\begin{qoptions}
\opt{A}{Magnetic field at origin at $t=0$ is $\left(-9\hat i+12\hat j\right)\times10^{-10}$T}
\opt{B}{Magnetic field at origin at $t=\dfrac{2\pi}{5}$ sec is $\left(-9\hat i+12\hat j\right)\times10^{-10}$T}
\opt{C}{Rate of change of magnetic field at $t=0$ at origin is $2.25\times10^{-9}$ T/s}
\opt{D}{Rate of change of magnetic field at $t=0$ at origin is zero}
\end{qoptions}
\qn{6} A uniform conducting rod PQ of mass 2 kg is hanged from roof using two non conducting strings of negligible mass length of rod is 2 m. An infinitely long current carrying wire is kept perpendicular to $xz$ plane and passes through $(1,1,1)$. 40 A current flow through wire in $+y$ direction. Tension in string connected at point P and Q are $T_1$ and $T_2$ respectively. If 20 A current is passed through rod PQ from P to Q then -
\qfig{p2_phy_s2_q06_1}
\begin{qoptions}
\opt{A}{Tension $T_1$ increases by $3.44\times10^{-5}$ N}
\opt{B}{Tension $T_1$ decreases by $8\times10^{-5}$ N}
\opt{C}{Tension $T_2$ increases by $3.44\times10^{-5}$ N}
\opt{D}{Tension $T_2$ decreases by $8\times10^{-5}$ N}
\end{qoptions}
\qn{7} In the figure shown a dipole is kept at the centre of a uniformly charged fixed ring. Charge on ring is $Q$ and dipole moment of dipole is $p$. Mass of dipole is $m$. dipole moment is along the axis of ring. Dipole is now released. Which of the following is/are true\,? (Neglect effect of gravity)
\qfig{p2_phy_s2_q07_1}
\begin{qoptions}
\opt{A}{Kinetic energy of dipole is maximum at a distance $\dfrac{R}{\sqrt2}$ from centre}
\opt{B}{Kinetic energy of dipole is maximum at distance $R$ from centre.}
\opt{C}{Maximum kinetic energy of dipole is $\dfrac{pQ}{6\sqrt3\,\pi\varepsilon_0R^2}$}
\opt{D}{Maximum kinetic energy of dipole is $\dfrac{pQ}{8\sqrt2\,\pi\varepsilon_0R^2}$}
\end{qoptions}
\sect{Section-I (iii)}
\qn{8} A uniform rod of mass $m$ and length $\ell$ has total thirteen holes of negligible size. All are equally spaced and are marked 1 to 13 from one end to another. Rod is made to oscillate about one of these holes. Match the hole number given in list-I to corresponding time period of oscillation given in list-II
\begin{center}\small\renewcommand{\arraystretch}{2.3}\setlength{\tabcolsep}{3pt}
\begin{tabular}{|l|l||c|l|}\hline
\multicolumn{2}{|c||}{List-I} & \multicolumn{2}{c|}{List-II}\\\hline
P. & Hole number-3 & 1. & $2\pi\sqrt{\dfrac{7\ell}{12g}}$\\\hline
Q. & Hole number-9 & 2. & $2\pi\sqrt{\dfrac{2\ell}{3g}}$\\\hline
R. & Hole number-4 & 3. & $2\pi\sqrt{\dfrac{4\ell}{3g}}$\\\hline
S. & Hole number-12 & 4. & $2\pi\sqrt{\dfrac{19\ell}{24g}}$\\\hline
 & & 5. & $2\pi\sqrt{\dfrac{37\ell}{60g}}$\\\hline
\end{tabular}
\end{center}
\begin{qoptions}
\opt{A}{$P\to1;\ Q\to3;\ R\to2;\ S\to4$}
\opt{B}{$P\to1;\ Q\to2;\ R\to3;\ S\to4$}
\opt{C}{$P\to1;\ Q\to2;\ R\to1;\ S\to5$}
\opt{D}{$P\to3;\ Q\to2;\ R\to4;\ S\to5$}
\end{qoptions}
\qn{9} A charged particle $\left(\dfrac qm=10\text{ c/kg}\right)$ is projected from origin with velocity $6\hat i+8\hat k$.
\qfig{p2_phy_s3_q09_1}
The magnetic field in uniform perpendicular to $x$ axis but varies as shown in graph with $x$ co-ordinate. Values of $x_0$ are given in List-I. Match the value of $y$ and $z$ coordinate of particle given in List-II when it crosses plane $x=4x_0$.
\begin{center}\small\renewcommand{\arraystretch}{1.3}\setlength{\tabcolsep}{3pt}
\begin{tabular}{|l|l||c|l|}\hline
\multicolumn{2}{|c||}{List-I} & \multicolumn{2}{c|}{List-II}\\\hline
P. & $x_0=12\pi$ m & 1. & $y=0,\ z=0$\\\hline
Q. & $x_0=6\pi$ m & 2. & $y=32$ m, $z=0$\\\hline
R. & $x_0=4\pi$ m & 3. & $y=64$ m, $z=0$\\\hline
S. & $x_0=3\pi$ m & 4. & $y=32$ m, $z=32$ m\\\hline
 & & 5. & $y=48$ m, $z=32\sqrt3$ m\\\hline
\end{tabular}
\end{center}
\begin{qoptions}
\opt{A}{$P\to2;\ Q\to3;\ R\to4;\ S\to5$}
\opt{B}{$P\to1;\ Q\to2;\ R\to3;\ S\to5$}
\opt{C}{$P\to1;\ Q\to3;\ R\to4;\ S\to5$}
\opt{D}{$P\to1;\ Q\to3;\ R\to5;\ S\to4$}
\end{qoptions}
\qn{10} A sphere of radius $R$ has uniform volume charge density $\rho$. A spherical cavity of radius $\dfrac R2$ is removed as shown. This cavity is filled with material having charge density $2\rho$. Match the point given in list-I to corresponding magnitude of strength of electric field given in list-II.
\qfig{p2_phy_s3_q10_1}
\begin{center}\small\renewcommand{\arraystretch}{2.3}\setlength{\tabcolsep}{3pt}
\begin{tabular}{|l|l||c|l|}\hline
\multicolumn{2}{|c||}{List-I} & \multicolumn{2}{c|}{List-II}\\\hline
P. & $A$ & 1. & $\dfrac{\rho R}{2\varepsilon_0}$\\\hline
Q. & $B$ & 2. & $\dfrac{\rho R}{6\varepsilon_0}$\\\hline
R. & $C$ & 3. & $\dfrac{3\rho R}{8\varepsilon_0}$\\\hline
S. & $D$ & 4. & $\dfrac{19\rho R}{54\varepsilon_0}$\\\hline
 & & 5. & $\dfrac{17\rho R}{54\varepsilon_0}$\\\hline
\end{tabular}
\end{center}
\begin{qoptions}
\opt{A}{$P\to3;\ Q\to1;\ R\to2;\ S\to4$}
\opt{B}{$P\to1;\ Q\to2;\ R\to2;\ S\to4$}
\opt{C}{$P\to3;\ Q\to2;\ R\to1;\ S\to5$}
\opt{D}{$P\to1;\ Q\to2;\ R\to3;\ S\to5$}
\end{qoptions}
\qn{11} Three point changes are kept at various places of a cube. Charge $Q$, $2Q$ and $3Q$ are at midpoint of side AB, centre of cube and corner of cube respectively. Match the surfaces given in list-I to the corresponding flux passing through it given in list-II.
\qfig{p2_phy_s3_q11_1}
\begin{center}\small\renewcommand{\arraystretch}{2.3}\setlength{\tabcolsep}{3pt}
\begin{tabular}{|l|l||c|l|}\hline
\multicolumn{2}{|c||}{List-I} & \multicolumn{2}{c|}{List-II}\\\hline
P. & Total flux through cube & 1. & $\dfrac{Q}{3\varepsilon_0}$\\\hline
Q. & Flux through surface ABCD & 2. & $\dfrac{11Q}{24\varepsilon_0}$\\\hline
R. & Flux through surface ABFE & 3. & $\dfrac{11Q}{12\varepsilon_0}$\\\hline
S. & Sum of flux through BCGF and EFGH & 4. & $\dfrac{25Q}{24\varepsilon_0}$\\\hline
 & & 5. & $\dfrac{21Q}{8\varepsilon_0}$\\\hline
\end{tabular}
\end{center}
\begin{qoptions}
\opt{A}{$P\to5;\ Q\to1;\ R\to3;\ S\to2$}
\opt{B}{$P\to5;\ Q\to1;\ R\to2;\ S\to3$}
\opt{C}{$P\to5;\ Q\to1;\ R\to2;\ S\to4$}
\opt{D}{$P\to4;\ Q\to1;\ R\to3;\ S\to2$}
\end{qoptions}
\sect{Section-II}
\qn{1} A current is flowing through a cuboidal wire along its length as shown. When a uniform magnetic field of 200 G is switched on perpendicular to front face, a potential difference of 0.5 $\mu$V is detected across its top and bottom surface. If number density of electron in cuboidal metal conductor is $p\times10^{18}/\text{cm}^3$ then $p=$ \_\_\_\_\_\_\_\_\_\_\_\_
\qfig{p2_phy_s4_q01_1}
\qn{2} A point charge $Q$ is kept at the centre of a hollow conducting sphere of inner radius $a$ and outer radius $2a$. Conducting sphere is further coated with a dielectric medium having dielectric constant 3. Inner radius of dielectric is $2a$ and outer radius is $3a$. Potential at distance of $\dfrac a2$ from point charge is $x\times\dfrac{kQ}{72a}$ then $x=$ \_\_\_\_\_\_\_\_\_\_\_
\qfig{p2_phy_s4_q02_1}
\qn{3} Sun subtends an angle of 10 m rad on earth.\\
Calculate average density of sun is SI unit in nearest integer.\\
{[}Mean density of earth $=5.5$ g/cm$^3$, 1 year $\approx3\times10^7$ sec, $\pi=3.14$]
\qn{4} Density of a spherical planet varies linearly from center to surface as $\rho=\rho_0\left(2-\dfrac{r^2}{R^2}\right)$, where $r$ is distance from centre. A particle of mass 1 kg is situated at centre. If potential of energy of particle is $-q\times10^6$ J then, $q=$ \_\_\_\_\_\_\_\_\_\_\_\_\_\\
(Given radius of planet $=2000$ km, $G=\dfrac{20}{3}\times10^{-11}\ \text{N}\times\text{m}^2/\text{kg}^2$, $\rho_0=\dfrac{10^4}{\pi}\ \text{kg/m}^3$)
\qn{5} In the given circuit, all capacitor have capacitance of 10$\mu$F and EMF of batteries is 10 V. All capacitors are initially uncharged. Calculate charge on capacitor marked P at steady state in $\mu$C.
\qfig{p2_phy_s4_q05_1}
\qn{6} A rectangular uniform plate of mass $m$ has length $2\ell$ and width $\ell$ as shown. It is connected to four identical springs of spring constant $K$ at its corners. It is rotated by small angle about an axis passing through centre of plate and perpendicular to side AB. It angular frequency of plate is $\sqrt{\dfrac{xk}{2m}}$ then $x=$\_\_\_\_\_\_
\qfig{p2_phy_s4_q06_1}
\qfullbegin
\partsec{2}{CHEMISTRY}
\sect{Section-I (i)}
\qn{1} Dry air is passed through various containers as shown
\qfig{p2_chem_s1_q01_1}
\qfullsplit
The change in the masses of P, S and D are 3 gram, 2 gram and 1 gram respectively. Determine RLVP for the solution
\begin{qoptions}
\opt{A}{$\dfrac13$}
\opt{B}{$\dfrac23$}
\opt{C}{$\dfrac15$}
\opt{D}{$\dfrac25$}
\end{qoptions}
\qfullend
\qn{2} Ammoniated complex of cobalt ions shows yellow colour, which then converts into brown on exposure to air. This brown complex is further treated with glycine in presence of ammonium hydroxide. Which of the following is /are INCORRECT for the described process?
\begin{qoptions}
\opt{A}{Oxidized ammoniated complex is diamagnetic.}
\opt{B}{Glycinate complex shows two isomers.}
\opt{C}{One of the isomers of glycinate complex has maximum 3 identical atoms in a plane.}
\opt{D}{All the isomers of glycinate complex are optically active}
\end{qoptions}
\qn{3} Which of the following doesn't have a `B$-$X' back bond. Where `X' is an atom covalently bonded to Boron?
\begin{qoptions}
\opt{A}{$\mathrm{BF_3}$}
\opt{B}{$\mathrm{B(OH)_3}$}
\opt{C}{$\mathrm{BCl_3}$}
\opt{D}{$\left[\mathrm{B(OH)_4}\right]^-$}
\end{qoptions}
\qn{4} A newly discovered Element `Ai' has two isotopes $\mathrm{Ai^{400}}$ \& $\mathrm{Ai^{500}}$. The mass percent of each of the isotopes is 50\%. Find the average molar mass of `Ai'
\begin{qoptions}
\opt{A}{450 g mol$^{-1}$}
\opt{B}{444.44 g mol$^{-1}$}
\opt{C}{555.55 g mol$^{-1}$}
\opt{D}{455.56 g mol$^{-1}$}
\end{qoptions}
\sect{Section-I (ii)}
\qn{5} NaCl(aq) solution is at $0^\circ$C \& 1 atm. Now, some ice at $0^\circ$C is added into it, then:
\begin{qoptions}
\opt{A}{Vapour pressure of solution increases}
\opt{B}{Osmotic pressure of solution increase}
\opt{C}{Some ice will remain unmelted}
\opt{D}{The freezing point will increase}
\end{qoptions}
\qn{6} Which of the following is/are correct:
\begin{qoptions}
\opt{A}{The 3$^{\text{rd}}$ group of the long form of the modern periodic table has the highest number of elements among all the groups.}
\opt{B}{S $>$ Se $>$ O is the correct order of electron affinity for the elements}
\opt{C}{Cr, Mo \& W, all the three have an exceptional electronic configuration.}
\opt{D}{Ba $>$ Be $>$ B is the correct order of the atomic size for the elements.}
\end{qoptions}
\qn{7} Consider the reactions:
\qfig{p2_chem_s2_q07_1}
The correct relation between $\mathrm{P_1}$, $\mathrm{P_2}$ and $\mathrm{P_3}$, the major products, is/are
\begin{qoptions}
\opt{A}{$\mathrm{P_1}$ \& $\mathrm{P_2}$ are positional isomers}
\opt{B}{$\mathrm{P_2}$ \& $\mathrm{P_3}$ are enantiomers}
\opt{C}{$\mathrm{P_3}$ \& $\mathrm{P_1}$ are identical}
\opt{D}{$\mathrm{P_1}$ \& $\mathrm{P_2}$ are identical}
\end{qoptions}
\sect{Section-I (iii)}
\qn{8} Match the Columns\\
\textbf{Column-I} (Reactants $\to$ Major Product)\par
P.\ \qfigi{p2_chem_s3_q08_1}\par
Q.\ \qfigi{p2_chem_s3_q08_2}\par
R.\ \qfigi{p2_chem_s3_q08_3}\par
S.\ \qfigi{p2_chem_s3_q08_4}\par
\begin{center}\small
\begin{tabular}{|c|l|}\hline
\multicolumn{2}{|c|}{\textbf{Column-II} (Appropriate reagent)}\\\hline
1. & NBS\\\hline
2. & $\mathrm{Br_2/CCl_4}$\\\hline
3. & HBr\\\hline
4. & $\mathrm{HBr/H_2O_2}$\\\hline
5. & Free radical intermediate involved\\\hline
\end{tabular}
\end{center}
\begin{qoptions}
\opt{A}{$P\to1,5;\ Q\to2,5;\ R\to3,5;\ S\to4,5$}
\opt{B}{$P\to4,5;\ Q\to1,5;\ R\to2,5;\ S\to3,5$}
\opt{C}{$P\to3;\ Q\to4,5;\ R\to1,5;\ S\to2$}
\opt{D}{$P\to2,5;\ Q\to3,5;\ R\to4,5;\ S\to1$}
\end{qoptions}
\qn{9} Match the Columns
\begin{center}\small
\begin{tabular}{|l|p{0.18\linewidth}|c|p{0.46\linewidth}|}\hline
\multicolumn{2}{|p{0.22\linewidth}|}{\centering\textbf{Column-I} Type of solid} & \multicolumn{2}{p{0.55\linewidth}|}{\centering\textbf{Column-II} Facts related to the solids assuming ideal radius ratio}\\\hline
P. & NaCl & 1. & Smaller ions present in the voids has cordination number greater than 4\\\hline
Q. & CsCl & 2. & The distance between the nearest voids containing the smaller ions is the same as the distance between any two nearest larger ions\\\hline
R. & ZnS (Zinc Blende) & 3. & 100\,\% of the voids of a particular type are occupied by the smaller ions.\\\hline
S. & $\mathrm{CaF_2}$ & 4. & 50\,\% of the voids of a particular type are occupied by the smaller ions.\\\hline
 & & 5. & Cordination number of the smaller ion is not more than the total number of formula units of the ionic compound in a single smallest unit cell.\\\hline
\end{tabular}
\end{center}
\begin{qoptions}
\opt{A}{$P\to1,2,3;\ Q\to1,2,3;\ R\to2,4,5;\ S\to3,5$}
\opt{B}{$P\to1,2,3,5;\ Q\to1,2,3,5;\ R\to2,4;\ S\to3$}
\opt{C}{$P\to1,3,5;\ Q\to1,3,5;\ R\to4,5;\ S\to2,3,5$}
\opt{D}{$P\to1,3;\ Q\to1,3;\ R\to2,3,5;\ S\to4,5$}
\end{qoptions}
\qn{10} Match the columns
\begin{center}\small
\begin{tabular}{|l|p{0.3\linewidth}|c|p{0.3\linewidth}|}\hline
\multicolumn{2}{|p{0.3\linewidth}|}{\centering\textbf{Column-I} Compounds} & \multicolumn{2}{p{0.4\linewidth}|}{\centering\textbf{Column-II} Characteristic properties of the compounds}\\\hline
P. & $\mathrm{Fe(NH_3)_6^{2+}}$ & 1. & Hybridization is $\mathrm{d^2sp^3}$\\\hline
Q. & $\mathrm{NiF_6^{2-}}$ & 2. & Spin free complex\\\hline
R & $\left[\mathrm{Co(H_2O)_6}\right]^{3+}$ & 3. & 18 electrons in valence shell\\\hline
S. & $\left[\mathrm{Pt\,(NH_3)_4Cl_2}\right]\mathrm{Cl_2}$ & 4. & $\mu_{\text{exp.}}\neq\mu_{\text{theor.}}$ ($\mu$ : magnetic moment)\\\hline
 & & 5. & Shape is perfectly octahedral\\\hline
\end{tabular}
\end{center}
\begin{qoptions}
\opt{A}{$P\to2,4;\ Q\to1,3,5;\ R\to1,3,5;\ S\to1,3$}
\opt{B}{$P\to2,3;\ Q\to1,2,3;\ R\to1,2;\ S\to2,4$}
\opt{C}{$P\to1,2;\ Q\to1,2,4;\ R\to1,2;\ S\to2,3$}
\opt{D}{$P\to1,2;\ Q\to1,2,5;\ R\to1,2;\ S\to1,3$}
\end{qoptions}
\qn{11} Match the column I to column-II\\
\textbf{Column-I} (Reactions)\par
P.\ \qfigi{p2_chem_s3_q11_1}\par
Q.\ \qfigi{p2_chem_s3_q11_2}\par
R.\ \qfigi{p2_chem_s3_q11_3}\par
S.\ \qfigi{p2_chem_s3_q11_4}\;$\xrightarrow[\text{(b) KOH(aq)}]{\text{(a) }\mathrm{HBr/H_2O_2}}$\par
\begin{center}\small
\begin{tabular}{|c|p{0.75\linewidth}|}\hline
\multicolumn{2}{|c|}{\textbf{Column-II} (Fact about the reaction)}\\\hline
1. & Addition of $\mathrm{H_2O}$ on the double bond to give Markownikoff's product\\\hline
2. & Addition of $\mathrm{H_2O}$ on the double bond to give Anti Markownikoff's product\\\hline
3. & Final product is the same as the final product of another reaction given in the column-I\\\hline
4. & Rearrangement of Carbocation involved\\\hline
5. & \qfigi{p2_chem_s3_q11_5}\\\hline
\end{tabular}
\end{center}
\begin{qoptions}
\opt{A}{$P\to1;\ Q\to2;\ R\to3;\ S\to4$}
\opt{B}{$P\to2;\ Q\to3;\ R\to4;\ S\to1$}
\opt{C}{$P\to3;\ Q\to4;\ R\to1;\ S\to2$}
\opt{D}{$P\to1;\ Q\to2;\ R\to3;\ S\to4$}
\end{qoptions}
\sect{Section-II}
\qn{1} 40 ml 1M $\mathrm{BaCl_2}$(aq) is mixed into a 60 ml solution of 1M $\mathrm{H_2SO_4}$(aq.). Where they react,\\
\centerline{$\mathrm{BaCl_2+H_2SO_4\longrightarrow BaSO_4+2HCl}$}
The molarity of $\mathrm{Ba^{2+}}$(aq) in the resultant solution is
\qn{2} Let, $\mathbf a$ be the total number of ligands which can act as didentate / polydentate ligands with identical donor sites.\\
DMG, en, Phen, Gly, acac, $\mathrm{SO_4^{2-}}$, pyridine, $\mathrm{PPh_3}$\\
And\\
$\mathbf b$ = (number of Geometrical isomers of $\left[\mathrm{Pt(H_2NCH(CH_3)COO^-)_2}\right]$ complex) + number of stereo isomers of $\left[\mathrm{Pt(gly)_3}\right]^+$ complex ion).\\
Then the value of $a+b$ is:
\qn{3} \qfig{p2_chem_s4_q03_1}
($\sim\!\!\sim$ $\Rightarrow$ Bond configuration not fixed, wedge $\Rightarrow$ bond configuration is fixed)\\
Find the total number of possible products (A)?
\qn{4} Consider a two chamber vessel. An SPM Separates the two chambers into equal volume as shown
\qfig{p2_chem_s4_q04_1}
A one liter aqueous solution $\mathbf X$, (density$=1.6$ g/ml) of a non volatile solute is filled in one of the chambers. Soon after that an equal volume of pure water (density$=1.0$ g/ml) is filled in the other chamber. Very quickly, the solution surface is now covered with a capillary of 50 cm$^2$ cross section. Pure water now starts flowing from the solvent towards the solution. At equilibrium condition the level of the solution in the capillary is 30.4 cm then find the osmotic pressure (in milli-atm) of the solution X. Assume $\Delta V_{\text{mix}}=0$\\[2pt]
Use: density of Hg $=13.4$ g/ml, atmospheric pressure $=76$ cm of Hg and $\dfrac{23.6}{176}=0.134$)\\
(round off to the nearest integer)
\qfig{p2_chem_s4_q04_2}
\qn{5} How many species involve sp mixing?\\
\centerline{$\mathrm{B_2,\ O_2,\ F_2,\ N_2,\ O_2^+,\ F_2^-,\ C_2,\ O_2}$}
\qn{6}~
\[\text{Smallest }2^\circ\text{ alcohol}\xrightarrow[\mathrm{H_3O^+}]{\mathrm{K_2Cr_2O_7}}A\xrightarrow[2.\,\mathrm{H_2O}]{1.\,\mathrm{EtMgBr}}B\]
\[B\xrightarrow{\mathrm{HBr}}C\xrightarrow[\mathrm{EtOH}]{\mathrm{EtONa}}D\xrightarrow[\mathrm{OH^-}]{\text{Cold }\mathrm{KMnO_4}}E\xrightarrow[\Delta]{\mathrm{H^+}}F\]
Molar mass of F is :
\partsec{3}{MATHEMATICS}
\sect{Section-I (i)}
\qn{1} If $\sin^{-1}x,\ \cos^{-1}x$ are inverse circular function. The principal values of above function are given below in their respective domain:\\
$\sin^{-1}x\in\left[\dfrac{3\pi}{2},\dfrac{5\pi}{2}\right],\ \cos^{-1}x\in[2\pi,3\pi]$\\
If $\sin^{-1}\sin(10)=a\pi+b$ (where $a,b\in I$) then $|a|+|b|$ is equal to :
\begin{qoptions}
\opt{A}{13}
\opt{B}{14}
\opt{C}{12}
\opt{D}{15}
\end{qoptions}
\qn{2} Let $x_1=\tan^{-1}\dfrac{2}{\sqrt3}>x_2>x_3\ \ldots\ldots$ are positive real numbers satisfying $\sin(x_{n+1}-x_n)+3^{\frac{1-2n}{2}}\sin x_n\sin x_{n+1}=0$ for $n\ge1$ then $\lim\limits_{n\to\infty}x_n$ equals :
\begin{qoptions}
\opt{A}{$-\dfrac\pi8$}
\opt{B}{$\dfrac\pi6$}
\opt{C}{$\dfrac\pi{10}$}
\opt{D}{$\dfrac\pi{12}$}
\end{qoptions}
\qn{3} If the quadratic equation $9a(b-c)x^2-3b(c-a)x+c(a-b)=0$ and $9x^2-3x-2m=0$ have a common root such that first equation has equal roots then number of values of $x\in[0,\pi]$ where $f(x)=\left[4\sin(4m-3)x-(4m+3)\right]$ is non-derivable is\\
(where $a,b,c$ and $m\in R$ and $[.]$ represents G.I.F.)
\begin{qoptions}
\opt{A}{7}
\opt{B}{8}
\opt{C}{9}
\opt{D}{10}
\end{qoptions}
\qn{4} Let $f:R\to R$ be a function defined by $f(x)=2x^3-21x^2+78x+24$. Number of integers in the solutions set of `$x$' satisfying the inequality $f\!\left(f\!\left(f(x)-2x^3\right)\right)\ge f\!\left(f\!\left(2x^3-f(x)\right)\right)$ is:
\begin{qoptions}
\opt{A}{3}
\opt{B}{4}
\opt{C}{5}
\opt{D}{6}
\end{qoptions}
\sect{Section-I (ii)}
\qn{5} Let $f(x)=\begin{cases}\sin\pi x; & 0\le x\le2\\ 2-x; & 2<x\le5\end{cases}$, $g(x)=\begin{cases}\max(|x|,\,1-|x|,\,|x-1|); & \text{If }x\in(0,1)\\ \min(|x|,\,1-|x|,\,|x-1|); & \text{other wise}\end{cases}$ then :
\begin{qoptions}
\opt{A}{$f\!\left(g\!\left(\dfrac14\right)\right)=\dfrac1{\sqrt2}$}
\opt{B}{$f(g(-1))=1$}
\opt{C}{$f(|g(-2)|)=0$}
\opt{D}{$f\!\left(g\!\left(\dfrac34\right)\right)=\dfrac1{\sqrt2}$}
\end{qoptions}
\qn{6} Let $f(x)=\begin{cases}\dfrac{e^{\frac1x}+e^{\frac2x}+e^{\frac3x}}{p\,e^{-2+\frac1x}+q\,e^{-1+\frac3x}}; & 0<x<\dfrac\pi2\\[3mm](\sin x+\cos x)^{\operatorname{cosec}x}; & -\dfrac\pi2<x<0\end{cases}$ and $\lim\limits_{x\to0}f(x)$ exists and finite then correct options can be:
\begin{qoptions}
\opt{A}{$p=2,\ q=1$}
\opt{B}{$p+q$ can be equal to 6}
\opt{C}{$p\in R,\ q=1$}
\opt{D}{$\log(p^2+q^2)\ge1$}
\end{qoptions}
\qn{7} Functions $f_r(x)$, $r=1,2,3,4$ are defined as $f_1(x)=\begin{cases}x^2; & 0\le x<1\\ 1; & x\ge1\\ 0; & \text{otherwise}\end{cases}$ \& $f_2(x)=f_1(-x),\ f_3(x)=f_2(-x),\ f_4(x)=f_3(-x)$ then :
\begin{qoptions}
\opt{A}{$f_1(x)f_2(x)=f_3(x)f_4(x)$}
\opt{B}{$f_3(x)f_4(x)=0\ \forall x\in R$}
\opt{C}{$f_2(-1/2)+f_3(1/2)=1/2$}
\opt{D}{$f_1(x)f_2(x)=0,\ \forall x\in R$}
\end{qoptions}
\sect{Section-I (iii)}
\qn{8} Match the list-I to list-II
\begin{center}\small
\begin{tabular}{|l|p{0.55\linewidth}|c|l|}\hline
\multicolumn{2}{|c|}{\textbf{List-I}} & \multicolumn{2}{c|}{\textbf{List-II}}\\\hline
P. & The number of values of $a$ such that $\lim\limits_{x\to a}\left[\sin^{-1}\dfrac{2x}{1+x^2}\right]$ does not exist (where $[.]$ G.I.F.) & 1. & 4\\\hline
Q. & If $f(x)=\{x+[\log_2(2023+x)]\}+\{x+[\log_2(2023+x^2)]\}+\ldots+\{x+[\log_2(2023+x^{10})]\}$, then $[f(e)]$ equals: (where $\{.\}\to$ Fractional Part function, $[.]\to$ Greatest Integer function.) & 2. & 2\\\hline
R. & The number of solution(s) of $f(f(x))=\dfrac x4$, where $f(x)=4x(1-x)$, $0\le x\le1$ & 3. & 7\\\hline
S. & If the number of element(s) in the range of $f(x)=\operatorname{sgn}\left(\left[|\sin x|+|\cos x|\right]\right)$ is $k$, then $6k$ is : (where $[.]\to$ G.I.F., $\operatorname{sgn}(x)$ denotes signum function) & 4. & 6\\\hline
 & & 5. & 5\\\hline
\end{tabular}
\end{center}
\begin{qoptions}
\opt{A}{$P\to4;\ Q\to1;\ R\to2;\ S\to3$}
\opt{B}{$P\to4;\ Q\to2;\ R\to1;\ S\to5$}
\opt{C}{$P\to5;\ Q\to3;\ R\to1;\ S\to4$}
\opt{D}{$P\to5;\ Q\to3;\ R\to2;\ S\to4$}
\end{qoptions}
\qn{9} Let $f(x+2y)=f(x)\cdot(f(y))^2\ \forall\ x,y\in R$.\\
If $f'(0)=\ln2$\\
(where $f(x)$ is a non negative function)
\begin{center}\small
\begin{tabular}{|l|p{0.5\linewidth}|c|l|}\hline
\multicolumn{2}{|c|}{\textbf{List-I}} & \multicolumn{2}{c|}{\textbf{List-II}}\\\hline
P. & The number of solution(s) of $f(x)+x=0$ is/are & 1. & 4\\\hline
Q. & The value of $f(\log_24)$ is & 2. & 3\\\hline
R. & The number of solution(s) of $f(x)-x^2=0$ is/are & 3. & 1\\\hline
S. & The number of solution(s) of $f(x)=2-x^2$ is/are & 4. & 2\\\hline
 & & 5. & 0\\\hline
\end{tabular}
\end{center}
\begin{qoptions}
\opt{A}{$P\to5;\ Q\to3;\ R\to4;\ S\to2$}
\opt{B}{$P\to3;\ Q\to1;\ R\to4;\ S\to4$}
\opt{C}{$P\to2;\ Q\to4;\ R\to1;\ S\to5$}
\opt{D}{$P\to3;\ Q\to1;\ R\to2;\ S\to4$}
\end{qoptions}
\qn{10} Math the column List - I and List -II
\begin{center}\small
\begin{tabular}{|l|p{0.55\linewidth}|c|l|}\hline
\multicolumn{2}{|c|}{\textbf{List-I}} & \multicolumn{2}{c|}{\textbf{List-II}}\\\hline
P. & If the shortest distance between the curve $y^2=x^3$ and $9x^2+9y^2-30y+16=0$ is `$d$', then $[d]$ is equal to (where $[.]\to$G.I.F.) & 1. & 9\\\hline
Q. & If the angle of intersection of the curve $y=\dfrac{x^2}{1+x^2}$ and $y=2^{-x}$ is $\tan^{-1}\left(\dfrac{2(1+\ln2)}{k-\ln2}\right)$, then $k$ is equal to : & 2. & 0\\\hline
R. & If $x+y-9=0$ is a normal to the parabola $y^2-2y-8x+17=0$ at $(a,b)$, then `$a+b-3$' is equal to : & 3. & 4\\\hline
S. & If $y=4x-5$ is tangent to the curve $y^2=Px^3+q$ at the point $(2,3)$, then $(p-q)$ is equal to : & 4. & 2\\\hline
 & & 5. & 6\\\hline
\end{tabular}
\end{center}
\begin{qoptions}
\opt{A}{$P\to2;\ Q\to3;\ R\to1;\ S\to4$}
\opt{B}{$P\to2;\ Q\to3;\ R\to5;\ S\to1$}
\opt{C}{$P\to2;\ Q\to4;\ R\to5;\ S\to1$}
\opt{D}{$P\to5;\ Q\to2;\ R\to1;\ S\to4$}
\end{qoptions}
\qn{11} Match the column.
\begin{center}\small
\begin{tabular}{|l|p{0.55\linewidth}|c|l|}\hline
\multicolumn{2}{|c|}{\textbf{List-I}} & \multicolumn{2}{c|}{\textbf{List-II}}\\\hline
P. & If $f(x)=2|x|+|x+2|-\big||2+x|-2|x|\big|$ and $f(x)$ is non differentiable at $m$ points, then `$2m$' is & 1. & 6\\\hline
Q. & If $f(x)=|x|+\{x\}+\operatorname{sgn}x+[-x]$ (where $[.]\to$ G.I.F. and $\{.\}\to$ F.P.F) and $f(x)$ is discontinuous at $n$ points in $(-3,3)$, then $n$ is : & 2. & 2\\\hline
R. & The number of local minimum of the function $f(x)=x^2+4\cos x+5$ is : & 3. & 8\\\hline
S. & If $f(x)=2|x|^3+3x^2-12|x|+1$ where $x\in[-1,2]$, then the greatest value of $f(x)$ in the given interval is : & 4. & 4\\\hline
 & & 5. & 5\\\hline
\end{tabular}
\end{center}
\begin{qoptions}
\opt{A}{$P\to3;\ Q\to4;\ R\to5;\ S\to1$}
\opt{B}{$P\to4;\ Q\to3;\ R\to2;\ S\to5$}
\opt{C}{$P\to3;\ Q\to4;\ R\to2;\ S\to5$}
\opt{D}{$P\to2;\ Q\to3;\ R\to4;\ S\to1$}
\end{qoptions}
\sect{Section-II}
\qn{1} Let $f(x)$ be a polynomial function satisfying $\dfrac{2x-5}{3}\le f(x)\le\dfrac{2x+7}{3}\ \ \forall x\in R$ and $\lim\limits_{x\to1}(f(x))^{\frac1{x-1}}=e^{\frac pq}$ (where $p,q\in N$) then least value of $p+q$ is
\qn{2} If $f(x)=\dfrac{\left(x+\frac1x\right)^6-\left(x^6+\frac1{x^6}\right)-2}{\left(x+\frac1x\right)^3+\left(x^3+\frac1{x^3}\right)}$, then absolute minimum value of $|f(x)|$ in its domain is
\qn{3} If $\lim\limits_{x\to0^+}\dfrac{\left(e\left(\frac1{1+x}\right)^{\frac1x}-1\right)}{x}=\ell$, then $4\ell$ is equal to
\qn{4} Let $f:[0,10]\to\{1,2,3\}$ the total number of function can be formed such that $f$ is discontinuous at all integral points only but right continuous at all odd integers as well as left continuous at all even integral points on its given domain are:
\qn{5} Let $f(x)$ be a differentiable function satisfying $f(y)-f(x)=\dfrac{y^{-y}}{x^{-x}}f\!\left(\dfrac{x^{-x}}{y^{-y}}\right)\ \forall\ x,y\in R^+$. If $f'(1)=1$, then the value of $\left|f(e)\cdot f\!\left(e^{-1}\right)\right|$ is
\qn{6} If the function $f(x)=2x^3-(8-a)x^2+\left(a^2+\dfrac{16}{9}\right)x-12$ has a local minima at some $x\in R^-$, then the number of integers in the range of $a$.
\end{multicols}
\end{document}
